% Auto-generated by paper/build_manuscript.py from paper/manuscript.md.
% Edit the markdown source and rerun the build script.
\documentclass[11pt]{article}

\usepackage[T1]{fontenc}
\usepackage[utf8]{inputenc}
\usepackage{lmodern}
\usepackage{microtype}
\usepackage[margin=1in]{geometry}
\usepackage{amsmath,amssymb,bm}
\usepackage{hyperref}
\usepackage{xurl}

\hypersetup{
  colorlinks=true,
  linkcolor=blue,
  urlcolor=blue
}

\providecommand{\tightlist}{%
  \setlength{\itemsep}{0pt}%
  \setlength{\parskip}{0pt}%
}

\title{Self-Mass Unobservability as a Restricted Effective Hypothesis: Internal Consistency, Weak-Field Handoff, and Provisional Joint Constraints}
\author{}
\date{2026-04-23}

\begin{document}
\maketitle

\begin{center}
\small
draft manuscript aligned to the repository state available in the current workspace\\
Repository: \texttt{lpaiu-cs/self-mass-unobservability}
\end{center}

\begin{abstract}
We study a proposed principle of \textbf{self-mass unobservability} and
ask a narrower question than its initial philosophical form suggests:
can any physically viable remnant of the idea survive as a testable
effective theory, and if so, what do the current local results actually
say? The work proceeds in staged form. First, we show that
self-subtraction by itself does not generate a new center-of-mass force:
through quadrupole order, the center-of-mass equation remains
independent of the self-subtraction parameter. Second, we show that a
literal internal-structure implementation is inconsistent with
self-gravitating compact objects: in Newtonian barotropes the
\(\lambda\to 1\) limit decompactifies the star rather than producing a
new equilibrium. These two results force a restricted interpretation in
which any anomaly can survive only in an effective passive-coupling
sector, a clock sector, or both.

On that restricted footing we analyze three observational branches. In
the weak field, a mock LLR program recovers injected linear signals at
the expected millimeter synodic scale while remaining nearly blind to
the quadratic term. In the strong field, a published-bound translator
for PSR J0337+1715 (inputs: Voisin et al.~2025) yields data-driven
conditional-slice constraints at the \(10^{-5}\) (\(\sigma_1\)) and
\(10^{-4}\) (\(\sigma_2\)) level, while the joint two-parameter
marginals remain an explicitly prior-box-conditional degeneracy ridge.
In the clock sector, a corrected leave-one-out \(\gamma\)-based audit
yields only a local amplitude constraint and does not by itself decide
the tied-versus-decoupled question. Combining the living local branches,
a provisional joint consistency scaffold shows that the current clock
audit is compatible with the strong-field free-fall posterior under the
tied relation, with the tied-versus-decoupled Bayes factor remaining a
prior-volume-dominated audit rather than a data-driven discriminant. The
current repository therefore supports a conservative conclusion:
self-mass unobservability is not viable as a literal structural
principle, but a restricted EFT survives as a coherent target for
further external weak-field and full-TOA tests.
\end{abstract}

\section{Introduction}\label{introduction}

The starting intuition of this project is that a gravitating body may
fail to ``observe'' its own mass in the same way it responds to
externally sourced gravity. Taken literally, that intuition is too vague
to count as a physical model. The first task is therefore not to assert
new physics, but to determine whether the idea can be made
mathematically consistent and observationally testable without
collapsing into contradiction or into already-known
equivalence-principle phenomenology.

The repository organizes that question into staged requests:

\begin{enumerate}
\def\labelenumi{\arabic{enumi}.}
\tightlist
\item
  analytic center-of-mass decoupling,
\item
  internal-structure consistency,
\item
  weak-field LLR mock sensitivity,
\item
  real-data LLR handoff and estimator audit,
\item
  strong-field free-fall translation for PSR J0337+1715,
\item
  clock-sector audit from binary-pulsar Einstein delay,
\item
  provisional joint consistency of the surviving local branches.
\end{enumerate}

The central conclusion of the staged program is already narrower than
the original philosophical slogan. The literal principle does
\textbf{not} survive as a structural law for self-gravitating matter.
What survives is, at most, a restricted effective description in which
self-gravity enters external coupling and/or redshift laws through
body-dependent sensitivity parameters.

\section{Restricted EFT formulation}\label{restricted-eft-formulation}

The surviving effective description used throughout the repository is

\[
\frac{m_{P,A}}{m_{I,A}} = 1 + \sigma_1 s_A + \sigma_2 s_A^2,
\qquad
s_A := \frac{E_{g,A}}{m_A c^2} > 0,
\]

together with a clock-sector deformation

\[
\frac{d\tau_A}{dt}
=
1-\frac{v_A^2}{2c^2}
-\left(1+\zeta_1 s_A+\zeta_2 s_A^2\right)\frac{U_{\rm ext}}{c^2}.
\]

Here:

\begin{itemize}
\tightlist
\item
  \(m_P\) is the passive gravitational mass,
\item
  \(m_I\) is the inertial mass,
\item
  \(s_A\) is a compactness-like self-gravity fraction,
\item
  \((\sigma_1,\sigma_2)\) parameterize free-fall anomalies,
\item
  \((\zeta_1,\zeta_2)\) parameterize clock anomalies.
\end{itemize}

Two model classes are then distinguished.

\subsection{Decoupled model}\label{decoupled-model}

The free-fall sector and clock sector are independent: \[
(\sigma_1,\sigma_2)\ \text{free},\qquad (\zeta_1,\zeta_2)\ \text{free}.
\]

\subsection{Tied model}\label{tied-model}

A metric-like relation is imposed: \[
\zeta_i = \sigma_i.
\]

The entire project eventually reduces to asking whether the available
branches prefer, require, or merely tolerate this tied identification.

\section{Internal consistency: what fails and what
survives}\label{internal-consistency-what-fails-and-what-survives}

\subsection{Center-of-mass motion does not acquire a new force from
self-subtraction}\label{center-of-mass-motion-does-not-acquire-a-new-force-from-self-subtraction}

For an extended body with density \(\rho(\boldsymbol\xi)\),
center-of-mass position \(\mathbf X\), and external potential
\(\Phi_{\rm ext}\), the staged derivation begins from

\[
L=\int d^3\xi\,\rho(\boldsymbol\xi)\left[
\frac12(\dot{\mathbf X}+\dot{\boldsymbol\xi})^2
-\Phi_{\rm ext}(\mathbf X+\boldsymbol\xi)
-(1-\lambda)\Phi_{\rm self}(\boldsymbol\xi)
\right].
\]

In the center-of-mass gauge and through quadrupole order, this becomes

\[
L
=
\frac12 M\dot{\mathbf X}^2
-M\Phi_{\rm ext}(\mathbf X)
-\frac12 Q^{ij}\partial_i\partial_j\Phi_{\rm ext}(\mathbf X)
+L_{\rm int}(\lambda)
+O(\ell^3\nabla^3\Phi_{\rm ext}).
\]

The crucial point is that \(\lambda\) appears only in the internal
sector \(L_{\rm int}\), not in the center-of-mass-dependent terms. The
resulting center-of-mass equation is therefore \(\lambda\)-independent
through quadrupole order. For spherical bodies, \(Q^{ij}=0\), so the
center of mass follows the ordinary test-body law at this order.

\textbf{Interpretation.} The simplest version of ``a body does not see
its own mass'' does \textbf{not} automatically generate a new free-fall
law. It merely changes internal bookkeeping unless additional coupling
rules are introduced.

\subsection{Literal structural self-unobservability is
inconsistent}\label{literal-structural-self-unobservability-is-inconsistent}

The second analytic test applies self-gravity suppression directly
inside hydrostatic structure equations. Writing \[
\alpha := 1-\lambda,
\] the Newtonian hydrostatic equations become \[
\frac{dP}{dr} = - \alpha G \frac{m(r)\rho(r)}{r^2},
\qquad
\frac{dm}{dr} = 4\pi r^2\rho(r).
\]

For the \(n=1\) polytrope \(P=K\rho^2\), one finds at fixed \(M\) and
\(K\): \[
R(\alpha)\propto \alpha^{-1/2},\qquad
\rho_c(\alpha)\propto \alpha^{3/2},\qquad
P_c(\alpha)\propto \alpha^3,\qquad
E_g(\alpha)\propto \alpha^{3/2}.
\]

Thus, as \(\alpha\to 0^+\) (equivalently \(\lambda\to 1\)), \[
R\to\infty,\qquad \rho_c\to 0,\qquad P_c\to 0,\qquad E_g\to 0.
\]

The object does not approach a new compact equilibrium. It
decompactifies into a zero-binding cloud.

\textbf{Interpretation.} Literal self-unobservability inside the star is
not a viable stellar-structure principle. Any surviving version of the
idea must therefore live in an effective coupling/redshift law for an
already-formed body, not in the structure equations themselves.

\section{Weak-field mock branch: what the LLR toy program really
shows}\label{weak-field-mock-branch-what-the-llr-toy-program-really-shows}

The weak-field branch was first developed as a controlled mock program
rather than a real LLR fit. The relevant combination is

\[
\Delta_{\rm SEP} = \sigma_1(s_E-s_M)+\sigma_2(s_E^2-s_M^2),
\]

and the standard weak-field synodic signal is modeled as \[
\delta r_{\cos D}\approx A_N\Delta_{\rm SEP},
\qquad
A_N = \frac{13.1\ {\rm m}}{s_E-s_M}.
\]

With the project's Earth and Moon self-gravity fractions, this becomes
approximately \[
\delta r_{\cos D}
\approx
(13.1\ {\rm m})\,\sigma_1
+
(6.32\times 10^{-9}\ {\rm m})\,\sigma_2.
\]

That coefficient already explains the central weak-field lesson:
LLR-like data are sensitive to the linear term and almost blind to the
quadratic term.

In the mock injection-recovery study:

\begin{itemize}
\tightlist
\item
  null injection gave \(\sigma_1=(1.82\pm 6.24)\times 10^{-5}\),
\item
  injected \(\sigma_1=5\times 10^{-4}\) was recovered as \[
  \sigma_1=(4.42\pm 0.64)\times 10^{-4},
  \]
\item
  the recovered synodic amplitude remained in the expected millimeter
  regime.
\end{itemize}

\textbf{Interpretation.} This branch is a \textbf{sensitivity study},
not a real-data likelihood. It justifies using weak-field LLR mainly as
a future probe of the linear free-fall coefficient \(\sigma_1\), while
leaving the quadratic direction \(\sigma_2\) primarily to strong-field
systems.

\section{Request 4: real-data LLR branch and why it stops
locally}\label{request-4-real-data-llr-branch-and-why-it-stops-locally}

The repository pushed Request 4 much further than a specification-only
placeholder. The data side and handoff side were both developed, but the
final local weak-field posterior still did \textbf{not} close in the
current workspace.

\subsection{What succeeded}\label{what-succeeded}

The following steps were completed:

\begin{enumerate}
\def\labelenumi{\arabic{enumi}.}
\tightlist
\item
  \textbf{Pinned APOLLO ingest.} A public APOLLO legacy-text
  normal-point release was downloaded, hashed, parsed, and summarized.
\item
  \textbf{APOLLO-only nuisance scaffold.} A narrow baseline
  nuisance/design layer was constructed.
\item
  \textbf{APOLLO-only baseline probe.} A first surrogate GR-like
  baseline fit reduced a geocentric \texttt{DE421} nominal residual by
  about \(12\times\) in RMS, but still left residuals at
  \(O(10^5\!-\!10^6\,{\rm m})\).
\item
  \textbf{Canonical CRD pivot.} Public lunar CRD monthly files were
  acquired and parsed into a larger observation ensemble.
\item
  \textbf{MLRS replay and interface gates.} The public
  \texttt{MLRS\ Lunar\ Code} bundle was modernized just enough to replay
  the two bundled public lunar sample workflows end-to-end, with exact
  \texttt{.frd} reproduction and only one-line \texttt{.npt} drift.
  Bounded recalc, state, Sun-driven, and weak-field-style local-seam
  injections also survived without widening into \texttt{PLEPH} edits or
  ephemeris regeneration.
\end{enumerate}

The strongest of these local gates was the weak-field
observation-operator test, which injected an externally computed state
correction \[
\delta \mathbf r_{\rm rel}(t)=A_N\Delta_{\rm SEP}\,\hat{\mathbf u}_{ES}(t),
\qquad
\delta \mathbf v_{\rm rel}(t)=A_N\Delta_{\rm SEP}\,\frac{d\hat{\mathbf u}_{ES}}{dt},
\] through the \texttt{jjreadnp} seam in a barycenter-consistent way.
Both bundled samples propagated this bounded weak-field correction
through the normal-point pipeline.

\subsection{What failed}\label{what-failed}

Despite that success, the final public weak-field path did not close.

When representative monthly public CRD full-rate cases were promoted
through
\texttt{split\ -\textgreater{}\ v2→v1\ -\textgreater{}\ normalization\ -\textgreater{}\ ldb\_crd},
the path could reach parser/recalc, but still did \textbf{not} produce a
usable residual-bearing \texttt{93} stream. The bounded follow-up search
concluded that the missing \[
\texttt{raw frd} \to \texttt{usable frcal / 93 stream}
\] promotion layer is not publicly exposed in the current workspace and
is plausibly station-internal, removed from the public bundle, or both.

\subsection{Consequence}\label{consequence}

The local stop rule for Request 4 is therefore \textbf{triggered}.

\begin{itemize}
\tightlist
\item
  \texttt{MLRS} remains valuable as a \textbf{bounded weak-field
  observation-operator laboratory}.
\item
  It does \textbf{not} close to a public-MLRS weak-field posterior path
  here.
\item
  The final weak-field posterior must be handed off to a more mature
  external LLR estimator.
\end{itemize}

A public-code scout identified \texttt{PEP} as the strongest external
estimator family, but the current host cannot build it due to pervasive
\texttt{REAL*10} usage and legacy numeric assumptions
(\texttt{-m128bit-long-double}). A bounded local legacy/x86 environment
probe was also negative on this host.

\textbf{Interpretation.} Request 4 is not dead; it is
\textbf{externalized}. Locally, the repository has established the
weak-field observation-operator side and the exact boundary where a
mature external estimator becomes necessary.

\section{Strong-field free-fall anchor: PSR J0337+1715 Phase
A}\label{strong-field-free-fall-anchor-psr-j03371715-phase-a}

The strong-field anchor currently implemented in-workspace is the Phase
A translation for PSR J0337+1715. Instead of re-running a three-body
timing analysis, it maps published bounds on \(\Delta\) into the EFT
parameters:

\[
\Delta_{\rm SMU}
=
\sigma_1 (s_{\rm NS}-s_{\rm WD})
+
\sigma_2 (s_{\rm NS}^2-s_{\rm WD}^2).
\]

The two input bounds are the most recent published limits, from Voisin
et al.~(2025), A\&A 693, A143 (arXiv:2411.10066): the optimistic input
is \(|\Delta|<1.5\times 10^{-6}\) (95\% CL, circum-ternary planet
timing-noise hypothesis) and the conservative input is
\(|\Delta|<2.3\times 10^{-6}\) (95\% CL, achromatic red-noise
hypothesis). Earlier limits (Archibald et al.~2018; Voisin et al.~2020)
are compatible with and superseded by these inputs.

With the wide EOS surrogate prior \(s_{\rm NS}\in[0.10,0.20]\), the
\textbf{data-driven headline numbers are the conditional slices}:

\begin{itemize}
\tightlist
\item
  optimistic, \(\sigma_2=0\): \[
  |\sigma_1|_{95}\sim 1.15\times 10^{-5},
  \]
\item
  optimistic, \(\sigma_1=0\): \[
  |\sigma_2|_{95}\sim 1.00\times 10^{-4},
  \]
\end{itemize}

with the conservative input giving \(1.75\times 10^{-5}\) and
\(1.53\times 10^{-4}\) respectively.

Because both \(\sigma_1\) and \(\sigma_2\) float simultaneously, the
joint posterior is ridge-like: the single published bound constrains one
EOS-smeared linear combination, leaving a slowly decaying degeneracy
ridge. The fully marginalized \(2D\) quantiles
(\(|\sigma_1|_{95}\sim 4.7\times 10^{-5}\),
\(|\sigma_2|_{95}\sim 3.57\times 10^{-4}\)) are therefore
\textbf{prior-box-conditional, not standalone constraints}: an explicit
box-scale diagnostic in the repository shows they grow essentially
linearly (\(\times 3.96/\times 4.04\)) when the flat prior box is
enlarged \(\times 4\), and they are identical for the optimistic and
conservative inputs, while the conditional slices stay flat under the
same test.

The project later established that Phase B is
\textbf{public-input-ready}: public TOA files, par files, and
\texttt{Nutimo} code releases are mirrored locally. But the current host
still lacks the runtime stack (\texttt{Boost}, \texttt{Tempo2},
\texttt{Minuit}/\texttt{root-config}, \texttt{cython}, OpenMP-compatible
compiler path), so a full local TOA refit is not closed here.

\textbf{Interpretation.} The strong-field branch is more than a
placeholder: it already supplies a usable free-fall anchor for joint
consistency, even though the full TOA realization remains an
external-environment task.

\section{Clock-sector audit: corrected but still
local}\label{clock-sector-audit-corrected-but-still-local}

The clock-sector branch was substantially repaired. The earlier toy
ansatz \[
\gamma_{\rm obs}=(1+\zeta_1 s+\zeta_2 s^2)\gamma_{\rm GR}
\] was replaced by the form implied by the EFT when only the
gravitational redshift piece is modified:

\[
\gamma_{\rm SMU}
=
\gamma_{\rm GR}
\left[
1+\frac{\eta_p}{1+X_c}
\right],
\qquad
\eta_p=\zeta_1 s_p+\zeta_2 s_p^2.
\]

A leave-one-out mass reconstruction was then enforced so that \(\gamma\)
was not reused to define the masses it was supposed to test.

\subsection{What the current sample
identifies}\label{what-the-current-sample-identifies}

The current clock likelihood does \textbf{not} separately identify
\(\zeta_1\) and \(\zeta_2\). It is dominated by a local basis near the
sampled self-gravity region: \[
\eta_* := \zeta_1 s_* + \zeta_2 s_*^2,\qquad
\kappa_* := \zeta_1 + 2s_* \zeta_2.
\]

A caveat applies to the two-system base stage: because the two base
systems sample nearly the same self-gravity fraction (spread
\(\sim 0.6\%\)), the \(\gamma\) data there constrain only the band
\(\eta(\bar s)\) (\(|\eta_*|_{95}=1.48\times 10^{-3}\)); the separate
\((\zeta_1,\zeta_2)\) marginals and the slope quantile of that grid
stage are set by the flat \(\zeta\) prior box, not by the data (an
explicit box-scale diagnostic in the repository shows them growing
linearly with the box, while the Fisher lever-arm audit of the same two
systems gives only \(|\kappa_*|_{95}\simeq 8.4\)).

After the corrected base fit and the later additions of
\texttt{B1913+16}, \texttt{J1141-6545}, and \texttt{J1906+0746}, the
branch still behaves as a local audit rather than a global slope
discriminator. At the combined stage (Fisher-style audit; the
\texttt{B1913+16} lever arm at \(s\simeq 0.16\) makes these slope
numbers approximately data-driven, unlike the base-stage grid
quantiles): \[
s_{\rm ref}=0.155833,\qquad
|\eta_*|_{95}=2.105\times 10^{-3},\qquad
|\kappa_*|_{95}=4.628\times 10^{-2}.
\]

The low-side bottleneck remains. Even with the two follow-up low-side
systems, the slope direction does not open to the \(10^{-2}\) scale.
Therefore the project's own stop rule demotes Request 6 from ``main
novelty engine'' to \textbf{support/local-audit branch}.

\textbf{Interpretation.} The clock branch is valid as a corrected
preliminary audit, but it is not the place where the
tied-versus-decoupled verdict will be decided.

\section{Provisional joint consistency
scaffold}\label{provisional-joint-consistency-scaffold}

Given the freezes on Requests 4, 5 Phase B, and 6, the highest-value
local result became a \textbf{joint scaffold} that combines the living
branches:

\begin{itemize}
\tightlist
\item
  Request 5 Phase A as the strong-field free-fall anchor,
\item
  Request 6 as the local clock audit,
\item
  Request 3 only as a scale sanity check, not as evidence.
\end{itemize}

\subsection{Models}\label{models}

\subsubsection{Decoupled model}\label{decoupled-model-1}

\[
\mathcal L_{\rm dec}
=
\mathcal L_{\rm J0337\,Phase\,A}(\sigma_1,\sigma_2)
\times
\mathcal L_{\rm clock}(\eta_*,\kappa_*).
\]

\subsubsection{Tied model}\label{tied-model-1}

\[
\zeta_i=\sigma_i,
\] with projection into the local clock basis \[
\eta_*=\sigma_1 s_*+\sigma_2 s_*^2,\qquad
\kappa_*=\sigma_1+2s_*\sigma_2.
\]

\subsection{Clock surrogate variants}\label{clock-surrogate-variants}

The clock branch enters as a Gaussian surrogate in \((\eta_*,\kappa_*)\)
built from the two-system Request 6 stage. Because the \(\kappa_*\)
width of that grid posterior is set by its \(\zeta\) prior box, the
scaffold sweeps two surrogate variants: \texttt{grid\_matched} (moments
taken directly from the grid; overstates the clock slope information)
and \texttt{eta\_only\_kappa\_free} (the reference variant: \(\eta_*\)
from the grid posterior, \(\kappa_*\) width from the Fisher lever-arm
audit, \(|\kappa_*|_{95}\simeq 8.42\)).

\subsection{Result}\label{result}

For the reference scenario
(\texttt{optimistic\ /\ wide\ /\ medium\ /\ eta\_only\_kappa\_free}),
the scaffold gives:

\begin{itemize}
\tightlist
\item
  pre-clock projected tension: \[
  0.591\ \text{Mahalanobis}
  \] (0.592 under \texttt{grid\_matched}; the tension is carried by the
  \(\eta_*\) direction, so the surrogate choice barely moves it),
\item
  tied posterior projection: \[
  |\eta_*|_{95}\approx 1.99\times 10^{-6}\ \text{(data-limited direction)},\qquad
  |\kappa_*|_{95}\approx 5.62\times 10^{-5}\ \text{(box-conditional, inherited from the Phase A ridge)},
  \]
\item
  standalone clock audit: \[
  |\eta_*|_{95}\approx 1.48\times 10^{-3}\ \text{(data-driven)},
  \] with the clock slope direction essentially unconstrained at the
  two-system stage (Fisher \(|\kappa_*|_{95}\simeq 8.42\)).
\end{itemize}

Thus the tied image of the J0337 Phase A posterior sits deep inside the
region allowed by the current local clock audit.

Across the full sweep (2 surrogate variants \(\times\) 2 bounds
\(\times\) 3 EOS priors \(\times\) 3 clock priors = 36 scenarios): \[
\log_{10} BF(\text{decoupled}/\text{tied}) \in [-2.41,-0.49],
\] (\([-2.41,-0.82]\) under \texttt{grid\_matched}, \([-1.49,-0.49]\)
under \texttt{eta\_only\_kappa\_free}). The variation is driven almost
entirely by the chosen clock-prior box: the clock surrogate likelihood
is nearly constant over the J0337 \(\sigma\) grid, so the J0337 factor
cancels in the evidence ratio and the Bayes factor reduces to an
Occam-volume ratio (only a handful of distinct values occur across all
36 scenarios). The projected tension and the J0337-averaged clock
likelihood remain nearly unchanged, the former structurally so, because
the symmetric J0337 posterior always projects to a near-zero mean.

\textbf{Interpretation.} The current local data do \textbf{not} reveal a
strong inconsistency between the strong-field free-fall anchor and the
clock local audit under the tied relation. Nor do they provide a
decisive empirical victory for the tied model. At present, model
comparison is a \textbf{prior-sensitivity audit}, not a data-driven
discriminant.

\section{What the repository now really
supports}\label{what-the-repository-now-really-supports}

The repository now supports the following conservative claims.

\subsection{Supported claims}\label{supported-claims}

\begin{enumerate}
\def\labelenumi{\arabic{enumi}.}
\item
  \textbf{Literal self-mass unobservability is not viable as a
  structural principle.}\\
  The internal-structure branch rules that out.
\item
  \textbf{A restricted EFT survives.}\\
  The only coherent remnant is an effective passive-coupling and/or
  clock-sector anomaly.
\item
  \textbf{Weak-field and strong-field roles are separated.}\\
  Weak-field LLR is mainly a future probe of the linear free-fall
  coefficient \(\sigma_1\), while strong-field systems are needed to
  constrain the quadratic direction \(\sigma_2\).
\item
  \textbf{The current clock sector is local and weak.}\\
  It is compatible with the strong-field anchor but does not decide the
  tied-versus-decoupled question.
\item
  \textbf{The decisive branches now lie outside the local host.}\\
  The final weak-field posterior needs an external mature LLR estimator,
  and the full \texttt{J0337} TOA refit needs a pre-provisioned runtime
  environment.
\end{enumerate}

\subsection{Unsupported claims}\label{unsupported-claims}

The repository does \textbf{not} currently justify any of the following
stronger claims:

\begin{itemize}
\tightlist
\item
  a final empirical verdict for tied versus decoupled physics,
\item
  a real weak-field LLR posterior,
\item
  a full self-consistent \texttt{J0337} TOA reanalysis,
\item
  a complete physical realization of Nordtvedt dynamics inside the local
  \texttt{MLRS} lab,
\item
  or standalone limits read off the joint \(2D\) marginals of the Phase
  A ridge or off the two-system clock-sector \((\zeta_1,\zeta_2)\) grid
  quantiles: the repository's box-scale diagnostics show those numbers
  track the flat prior boxes, so only the conditional slices (Phase A)
  and the local band \(\eta(\bar s)\) (clock sector) are data-driven at
  this stage.
\end{itemize}

These boundaries are not a weakness of presentation; they are the result
of explicit stop rules and bounded-search artifacts.

\section{Discussion}\label{discussion}

The main scientific value of the project is not that it already finds
new physics. Its value is that it \textbf{clarifies the boundary}
between what survives and what does not.

The most important conceptual transition is from a philosophical
statement---

\begin{quote}
``a body does not observe its own mass''
\end{quote}

---to a restricted and testable EFT. The project demonstrates that this
transition is not optional. Without it, the proposal either does nothing
to center-of-mass motion or destroys self-gravitating structure.

Methodologically, the repository also establishes a second boundary: the
distinction between \textbf{bounded local laboratory work} and
\textbf{final inference machinery}. The Request 4 \texttt{MLRS} path is
the clearest example. It shows that a local weak-field observation
operator can be made to live inside a bounded code seam, but it also
shows exactly where the public pathway stops and where a mature external
estimator becomes necessary.

Finally, the joint scaffold reveals the present scientific picture at
the local level: the strong-field free-fall anchor carries most of the
scale information, while the clock branch remains broad and local. The
immediate implication is not that the tied model is validated, but that
the current local evidence does \textbf{not} expose a sharp
contradiction with it.

\section{Conclusion}\label{conclusion}

The repository now supports a coherent, bounded, and publishable
intermediate conclusion.

\begin{enumerate}
\def\labelenumi{\arabic{enumi}.}
\tightlist
\item
  \textbf{Self-mass unobservability fails as a literal structural law.}
\item
  \textbf{A restricted EFT survives} in passive-coupling and/or clock
  sectors.
\item
  \textbf{Weak-field mock data} support the expected hierarchy:
  \(\sigma_1\) can be seen, \(\sigma_2\) is weak in LLR.
\item
  \textbf{The strong-field branch} already provides a usable free-fall
  anchor through the J0337 Phase A translator.
\item
  \textbf{The clock branch} remains a local audit and does not currently
  force decoupling.
\item
  \textbf{The provisional joint scaffold} shows compatibility, not
  contradiction, between the strong-field free-fall anchor and the
  present clock audit under the tied relation.
\item
  \textbf{The final verdict remains externalized} to a mature weak-field
  LLR estimator and a full \texttt{J0337} TOA branch.
\end{enumerate}

That is a narrower result than a discovery claim, but it is also a more
reliable one. The project has converted an initially vague idea into a
restricted EFT with explicit analytic boundaries, reproducible stop
rules, and a clear hand-off structure for the branches that require
heavier inference infrastructure.

\section{Data, code, and reproducibility
note}\label{data-code-and-reproducibility-note}

The repository is public and organized around staged
\texttt{REQUEST*.md} memos, executable \texttt{request*.py} scripts, and
checked-in generated artifacts. The current public \texttt{README}
explicitly presents the repository as a staged research workspace and
states that it does \textbf{not} yet contain a final weak-field
posterior for Request 4, that Request 5 Phase A is implemented while
Phase B stops at public-input/build-feasibility gates, that Request 6 is
a local clock-sector audit, and that Request 7 is a provisional joint
consistency scaffold. This manuscript is written to that
repository-defined scope and deliberately avoids stronger claims.

Two reproducibility upgrades accompany this revision: the Request 5 and
Request 6 grid stages now ship explicit \texttt{box\_scale\_check}
diagnostics that rerun each posterior with the flat prior box enlarged
\(\times 2\) and \(\times 4\) (separating data-limited from box-limited
quantiles, as used throughout Sections 6--8), and all Monte-Carlo stages
now use deterministic seeding (the earlier per-source offsets used
Python's process-salted \texttt{hash()}, which made checked-in outputs
non-reproducible). A \texttt{references.bib} file with the citation
metadata is included alongside the manuscript source.

\section{Selected references (draft; convert to BibTeX before
submission)}\label{selected-references-draft-convert-to-bibtex-before-submission}

\begin{enumerate}
\def\labelenumi{\arabic{enumi}.}
\tightlist
\item
  Chandler, J. F., Battat, J. B. R., Murphy, T. W., Reardon, D.,
  Reasenberg, R. D., and Shapiro, I. I. (2021). \emph{The Planetary
  Ephemeris Program: Capability, Comparison, and Open Source
  Availability}. arXiv:2103.16745.
\item
  APOLLO normal-point public archive. University of California San Diego
  APOLLO normal points page.
\item
  ILRS format documentation and CRD manuals, including the CRD overview
  and format manuals.
\item
  Voisin, G. et al.~(2020). \emph{An improved test of the strong
  equivalence principle with the pulsar in a triple star system}. A\&A
  638, A24. arXiv:2005.01388. Public data/code release for the PSR
  J0337+1715 SEP analysis.
\item
  Voisin, G. et al.~(2025). \emph{Explanation of the exceptionally
  strong timing noise of PSR J0337+1715 by a circum-ternary planet and
  consequences for gravity tests}. A\&A 693, A143. arXiv:2411.10066.
  Source of the \(|\Delta|<1.5\times 10^{-6}\) (planet hypothesis) and
  \(|\Delta|<2.3\times 10^{-6}\) (red-noise hypothesis) 95\% inputs used
  in Phase A.
\item
  Fonseca, E. et al.~(2014). Timing and post-Keplerian analysis for PSR
  B1534+12. arXiv:1402.4836.
\item
  Kramer, M. et al.~(2021). Double Pulsar timing analysis.
  arXiv:2112.06795.
\item
  Kramer, M. et al.~(2006). Double Pulsar mass-ratio result, as used in
  the later Request 6 reconstruction.
\item
  Weisberg, J. M., Nice, D. J., and Taylor, J. H. (2010). PSR B1913+16
  timing analysis. arXiv:1011.0718.
\item
  Bhat, N. D. R., Bailes, M., and Verbiest, J. P. W. (2008). PSR
  J1141-6545 timing inputs. arXiv:0804.0956.
\item
  Venkatraman Krishnan, V. et al.~(2020). Spin-orbit / $\dot{x}$
  structure in PSR J1141-6545. arXiv:2001.11405.
\item
  Vleeschower, L. et al.~(2026). \emph{Long-term timing of the
  relativistic binary PSR J1906+0746}, with explicit \$
  \dot{x}-\gamma \$ correlation. arXiv:2602.05947.
\item
  Cameron, A. D. et al.~(2017). PSR J1757-1854 timing constraints.
  arXiv:1711.07697.
\item
  Ferdman, R. D. et al.~(2014). PSR J1756-2251 timing constraints.
  arXiv:1406.5507.
\item
  Ferdman, R. D. et al.~(2020). PSR J1913+1102 timing constraints.
  arXiv:2007.04175.
\item
  Archibald, A. M. et al.~(2018). \emph{Universality of free fall from
  the orbital motion of a pulsar in a stellar triple system}. Nature
  559, 73. arXiv:1807.02059.
\end{enumerate}

\end{document}
